\documentclass[a4paper,12pt]{article}
\usepackage[utf8]{inputenc}
\usepackage[T1]{fontenc}
\usepackage[french]{babel}
\usepackage{geometry}
\geometry{margin=2.5cm}

\title{\textbf{Cours d'Histoire - Culture Générale}}
\author{Mignard Mael}
\date{\today}

\begin{document}

\maketitle

\section{La Préhistoire (\textit{-2,5 millions d'années / -3000 })}
Il s'agit de la plus longue période de l'histoire humaine, s'étendant de l'apparition des premiers hominidés il y a environ 2,5 millions d'années jusqu'à l'invention de l'écriture vers 3000 av. J.-C. La Préhistoire est divisée en trois grandes périodes : le Paléolithique, le Mésolithique et le Néolithique.
\subsection{L'apparition de l'homme}
Les premiere especes humaine apparaissent en Afrique, l'etre humain developpe plusieur caracteristique fondamentaless
\begin{itemize}
    \item  La creation d'outils en pierre pour chasser et se défendre.
    \item La maîtrise du feu pour cuisiner et se réchauffer.
    \item Le développement du langage pour communiquer et transmettre des connaissances.
    \item se deplacer en groupe (famille, tribu) pour mieux se proteger et partager les ressources.
\end{itemize}
\subsection{la revolution neolitique}
La revolution neolitique marque une transition majeure dans l'histoire humaine, caractérisée par le passage d'un mode de vie nomade de chasseurs-cueilleurs à un mode de vie sédentaire basé sur l'agriculture et l'élevage. Cette transition a eu lieu il y a environ 10 000 ans et a conduit à la formation des premières civilisations.

\section{Antiquité (\textit{-3000 / 476} )}
L'antiquité debute avec l'invention de l'écriture et correspond a l'emergence des premiere grandes civilisations.

Les société deviennent plus complexe, et les premieres grande villes apparraissent.

\subsection{Les grandes Civilsations.}
\begin{itemize}
    \item \textbf{Rome}: L'empire romain a creer un veritable reseau de routes, le droit et organise un vaste territoire autour de la mediterranee. \\
    La lange Latine, dominante de l'administration et de la culture.
\end{itemize}

\section{Moyen-Age (\textit{476 / 1492} )}
Souvent considéré a tord comment une periode de stagnation. En réalité il s'agit d'une epoque de transformation importante.

\subsection{}

\end{document}

